\VerdierPhysicalPageAnchor{0011}
\begin{description}
\VerdierSemanticAnchor{front-publications-12}
\item[{[12]}] \emph{Théorème de dualité pour la cohomologie des espaces
localement compacts}, Dualité de Poincaré, Séminaire Heidelberg-Strasbourg
1966-67, exp. 4, Publ. I.R.M.A. Strasbourg, n° 3, pp. 1-36 (1969).

\VerdierSemanticAnchor{front-publications-13}
\item[{[13]}] \emph{Théorème du noyau}, Dualité de Poincaré, Séminaire
Heidelberg-Strasbourg 1966-67, exp. 9, Publ. I.R.M.A. Strasbourg, n° 3,
pp. 1-52 (1969).

\VerdierSemanticAnchor{front-publications-14}
\item[{[14]}] \emph{On a theorem of Wilder}, Applications of categorical
algebra, Proc. Symp. Pure Math. 17, pp. 184-191, A.M.S. (1970).

\VerdierSemanticAnchor{front-publications-15}
\item[{[15]}] \emph{Ein einfacher Beweis des Kohärenzsatzes von Grauert}
(en collaboration avec R. Kiehl), Math. Ann. 195, pp. 24-50,
Springer (1971).

\VerdierSemanticAnchor{front-publications-16}
\item[{[16]}] \emph{Topologie sur les espaces de cohomologie d'un complexe de
faisceaux analytiques à cohomologie cohérente}, Bull. Soc. Math. France 99,
pp. 337-343, Gauthier-Villars (1971).

\VerdierSemanticAnchor{front-publications-17}
\item[{[17]}] \emph{Dualité relative en géométrie analytique complexe}
(en collaboration avec J.-P. Ramis et G. Ruget), Inv. Math. 13,
pp. 261-283, Springer (1971).

\VerdierSemanticAnchor{front-publications-18}
\item[{[18]}] S.G.A. 4 (1963-64) “Théorie des topos et cohomologie étale des
schémas” : \emph{Préfaisceaux} (en collaboration avec A. Grothendieck)
pp. 1-184 ; \emph{Topologies et faisceaux}, pp. 219-263 ;
\emph{Fonctorialité des catégories de faisceaux}, pp. 265-297 ;
\emph{Topos} (en collaboration avec A. Grothendieck), pp. 299-519,
Lecture Notes in Mathematics 269, Springer (1972).

\VerdierSemanticAnchor{front-publications-19}
\item[{[19]}] S.G.A. 4 (1963-64) “Théorie des topos et cohomologie étale des
schémas” : \emph{Cohomologie dans les topos}, pp. 1-62 ;
\emph{Conditions de finitude. Topos et sites fibrés. Applications aux questions
de passage à la limite} (en collaboration avec A. Grothendieck), pp. 163-336,
Lecture Notes in Mathematics 270, Springer (1972).

\VerdierSemanticAnchor{front-publications-20}
\item[{[20]}] \emph{Caractéristique d'Euler-Poincaré}, Bull. Soc. Math.
France 101, pp. 441-445, Gauthier-Villars (1973).

\VerdierSemanticAnchor{front-publications-21}
\item[{[21]}] \emph{Le théorème de Le Potier}, Séminaire de géométrie
analytique 1972-73, Différents aspects de la positivité, Astérisque 17,
pp. 68-78, S.M.F. (1974).

\VerdierSemanticAnchor{front-publications-22}
\item[{[22]}] \emph{Indépendance par rapport à $l$ des polynômes
caractéristiques des endomorphismes de Frobenius de la cohomologie $l$-adique}
(d'après P. Deligne), Séminaire Bourbaki 1972-73, n° 423,
Lecture Notes in Mathematics 383, pp. 98-115, Springer (1974).

\VerdierSemanticAnchor{front-publications-23}
\item[{[23]}] \emph{Quelques problèmes d'intersection en géométrie
riemannienne} (en collaboration avec P. Marry), J. Diff. Geom. 12,
pp. 345-376, A.M.S. (1977).
\end{description}

